\section{Related Work}

\paragraph{EEG representations and benchmarks.}
REVE was designed to handle heterogeneous EEG recordings
\citep{elouahidi2025reve}. NeuralBench provides a framework for comparing
neuroimaging models across datasets and tasks \citep{banville2026neuralbench}.
We follow its age-prediction context and training-budget conventions, but train
only the prediction heads. MIPDB is not an official NeuralBench score.

Recent studies illustrate the importance of this distinction. STST-JEPA reports
strong age prediction with an attentive probe of frozen embeddings in its own
training and evaluation setting \citep{segal2026ststjepa}. NeuroAtlas finds that
model rankings vary across EEG domains \citep{kontras2026neuroatlas}, while
comparisons of EEG feature extractors distinguish the usefulness of frozen
features from performance after task-specific adaptation
\citep{turgut2026feature}. Our study isolates head choice within one encoder and
one development protocol, then tests transfer across cohorts.

\paragraph{Probing and age prediction.}
Linear probes use a deliberately restricted predictor to examine what a
representation makes accessible \citep{alain2016probes}. Richer heads ask a
broader question: whether additional computation can recover useful information
from those features. Comparing them therefore requires matching data exposure
and training conditions.

EEG age-prediction studies have examined sleep, childhood development, and
validation across datasets \citep{sun2019brainage,vandenbosch2019age,
iyer2024growthchart,tveitstol2026robustness}. Here chronological age is a
prediction target for studying representations, rather than a proposed clinical
biomarker.

\paragraph{Uncertainty in external evaluation.}
External neuroimaging studies face uncertainty from both training and evaluation
samples \citep{rosenblatt2024power}. Resampling methods can account for multiple
sources of variation \citep{saravanan2020hierarchical}. We resample seeds and
participants while preserving paired predictions, and separately test the
seed-level differences on the observed cohort. Holm correction accounts for
the three primary candidate comparisons \citep{holm1979simple}.
